\documentclass[11pt,a4paper]{article}
\usepackage[T1]{fontenc}
\usepackage[utf8]{inputenc}
\usepackage[spanish,es-noshorthands]{babel}
\usepackage[margin=2.5cm]{geometry}
\usepackage{enumitem}
\setlist{itemsep=2pt}

\title{Calibración del Citiroc: programa de festejos}
\author{JJ}
\date{\today}

\begin{document}
\maketitle

\section*{Aperitivo}

El programa de festejos pide calibrar la electrónica + SiPMs. Usamos la tarjeta CAEN A5202 (dos Citiroc~1A, 64 canales), leída a través del concentrador DT5215 con el programa de
adquisición de datos (DAQ) de JMB. Este aperitivo explica
solo el orden de las medidas, luego hay un menú con instrucciones para cada estudio.  

\section*{Peak stretch vs Track \& hold}

El Citiroc guarda la salida del slow shaper de cada canal de una de dos
maneras. En {\bf peak stretch} sigue la señal entre el trigger y el hold y
guarda su máximo; el \texttt{HoldDelay} fija la duración de esa ventana. En
{\bf track \& hold} guarda el valor en el instante del hold; el
\texttt{HoldDelay} fija entonces en qué punto del pulso se toma la muestra.
El DAQ de JMB trabaja a día de hoy siempre en peak stretch, que,
en todo caso, es el modo con el que se hacen las medidas de
energía.

\section*{Primera ronda de medidas: peak stretch}

\begin{enumerate}
\item {\bf Pedestal y ruido} (menú 1). Con los SiPM por debajo de la tensión
  de ruptura (breakdown voltage, BV), queremos calcular la media y la anchura del pedestal de
  cada canal en función del shaping time y del \texttt{HoldDelay}.
\item {\bf Hold delay scan} (menú 2). Con los SiPM por encima del BV, queremos calcular el
  espectro de los pulsos oscuros (Dark Pulses, DP) frente al \texttt{HoldDelay}. En peak
  stretch el valor guardado crece mientras la ventana no llega al máximo del
  pulso y después se estabiliza; el scan da el \texttt{HoldDelay} mínimo que
  recoge el máximo y cuánto sube el pedestal al alargar la ventana.
\item {\bf Calibración con los DP} (menú 3). Los espectros de
  uno, dos y tres fotoelectrones (single photoelectron, SPE) de cada canal,
  a varios shaping times y overvoltages (tensión por encima del BV), que dan
  la ganancia de cada canal en cuentas de ADC por fotoelectrón, y la tasa de
  trigger frente al umbral, que da el crosstalk.
\end{enumerate}

\noindent Empezamos por peak stretch por dos razones: es lo único que el DAQ
permite hoy, y es el modo en que se tomarán los datos con cristal, así que
esta calibración es la que se usará entonces. Además, estos datos nos
sirven ya para comparar con la simulación.

\section*{La forma del pulso: osciloscopio}

La fracción de la señal que el chip guarda depende de la forma del pulso a
la salida del slow shaper, y esa forma es la entrada de la simulación que
más influencia tiene en la posición y la anchura del pico de 511~keV. La A5202 saca a
un conector la señal analógica de un canal elegido (la sonda analógica). Puede medirse con el osciloscopio (es fácil que ya lo hayas hecho).

\begin{itemize}
\item {\bf Acción:}
  \begin{itemize}
  \item SiPM a 5~V de overvoltage, a oscuras.
  \item En la pestaña {\bf Acquisition}: \texttt{AnalogProbe0} = SLOW\_HG.
  \item \texttt{ProbeChannel0} = un canal entre 0 y 31.
  \item \texttt{DigitalProbe0} = HOLD.
  \item \texttt{T0\_Out} = DPROBE.
  \item Trigger PTRG.
  \item Conectar la señal analógica, que sale por el conector MCX Probe~1
    de la tarjeta (cable RG174 con clavija MCX macho), a un canal del
    osciloscopio, con la entrada a 1~M$\Omega$.
  \item Conectar la señal digital, que sale por el LEMO T0-OUT del panel
    frontal (niveles LVTTL), a otro canal, también a 1~M$\Omega$.
  \item Disparar el osciloscopio en la señal analógica, a la mitad de la
    altura de un DP de un fotoelectrón.
  \item Lanzar un run, para que la tarjeta genere el hold.
  \end{itemize}
\item {\bf Resultado:} los DP de un fotoelectrón en la pantalla y
  el instante del hold.
\item {\bf Acción:}
  \begin{itemize}
  \item Ventana del osciloscopio de unos 500~ns, que incluya la caída
    completa del DP.
  \item Para \texttt{HG\_ShapingTime} = 87,5, 50 y 25~ns: activar el
    promedio del osciloscopio (average) sobre unos 1000 disparos y guardar
    la waveform promediada en un fichero CSV, desde el menú de guardado del
    osciloscopio.
  \item En los mismos tres ajustes: con el promedio desactivado, adquirir en
    modo single y guardar una decena de DP sueltos, un CSV por DP.
  \item En cada fichero, guardar los dos canales: la señal analógica y la
    del hold.
  \item Repetir a 87,5~ns con \texttt{AnalogProbe0} = SLOW\_LG.
  \end{itemize}
\item {\bf Resultado:} ficheros CSV de waveforms por cada ajuste,
  con el shaping y la sonda en el nombre.
\item {\bf Precaución:} al terminar, volver a poner \texttt{AnalogProbe0} =
  OFF y \texttt{T0\_Out} a su valor de partida: con la sonda conectada la
  ganancia del canal cambia ligeramente.
\end{itemize}

\section*{Track \& hold, más adelante}

Track \& hold no está disponible todavía en el DAQ. 

\section*{Qué necesita el cocinero}

\begin{itemize}
\item De cada run, los ficheros \texttt{.h5} que indica cada menú, y una
  línea en el logbook con el número de run, la nota y el paso a los que corresponde, y cualquier incidencia.
\item {\bf El logbook:} abrir una hoja de Google, compartirla conmigo y
  escribir en ella una fila por cada run del DAQ y por cada fichero del
  osciloscopio, con estas columnas:
  \begin{center}
  \small
  \begin{tabular}{ll}
  \hline
  columna & contenido \\
  \hline
  fecha y hora & al empezar el run o al guardar el fichero \\
  run o fichero & número de run del DAQ, o nombre del CSV \\
  menú y paso & por ejemplo, 1.2 (menú 1, paso 2) u osciloscopio \\
  \texttt{HG\_ShapingTime} & en ns \\
  \texttt{LG\_ShapingTime} & en ns \\
  \texttt{HoldDelay} & en ns \\
  \texttt{HV\_Vbias} & en V \\
  temperatura & la de la tarjeta, en $^\circ$C \\
  eventos & número de eventos del run \\
  incidencias & cualquier cosa rara; vacío si no hay \\
  \hline
  \end{tabular}
  \end{center}
\item Los ficheros CSV de waveforms del osciloscopio, con el canal, la
  tensión y las escalas de tiempo y de amplitud anotadas.
\item Si algo no se comporta como espera el menú, para y grita. 
\end{itemize}

\end{document}
